\section{本章小结}
\label{sec:fourier-summary}

本章开头问：用复指数为函数换上频率坐标，对哪些函数有意义，在什么意义下能还原，运算规则在什么条件下成立。

周期函数的回答是Fourier级数。复指数$e^{\mathrm ikx}$在归一化内积下标准正交，Fejér核的非负平均证明了它们完备，于是每个$L^2$周期函数都有范数收敛的展开，Parseval关系是无限维的勾股定理。这一展开描述的是几乎处处等价类，逐点与一致收敛要另加条件：分段$C^1$函数的级数在每点收敛到左右极限的平均，$C^2$函数的系数按$k^{-2}$衰减，级数一致收敛；跳跃函数的部分和不可能一致收敛，跳跃附近出现约$9\%$的Gibbs超调。

全空间上的函数没有周期，下标连成了$\mathbb R^d$，而平面波不再属于所讨论的空间，类比开始失效，每一步都要直接证明。对$L^1$函数，Fourier变换在每个频率处都有确定值，是有界、连续、在无穷远处趋于零的函数。变换的存在与反演是两个问题：$\widehat f$也可积时，借助Gauss核的恒等逼近，反演积分给出$f$的连续代表，与$f$几乎处处相等；方框这样的跳跃函数变换不可积，只能在局部正则条件下用对称截断逐点反演，跳跃点处仍得到平均值。

能量关系属于$L^2$。截取说明$L^1\cap L^2$在$L^2$中稠密，Plancherel恒等式先在交集上成立，再由稠密延拓把变换唯一地推广为$L^2$上的等距。所得的变换是双射，满足$\mathcal F\mathcal F=R$，保持内积，性质与酉矩阵相同，但频域元素只是等价类，不能一律在单个频率处求值。同一个反演公式，在三种情形下分别表示绝对收敛积分的几乎处处等式、对称截断的逐点极限与$L^2$范数极限。

运算规则说明频率坐标使若干算子对角化。平移只改相位，调制平移频率；线性变量替换$\bm A$使频率按$\bm A^{-\top}$变换并带来因子$\lvert\det\bm A\rvert^{-1}$；弱导数可积时，求导变成乘以$2\pi\mathrm i\xi_j$，光滑性因而对应高频衰减；卷积在$L^1$中变成逐点乘积，推广到$L^1\times L^2$与$L^2\times L^2$时，输出所在的空间与等式的意义都要重新核对。

离散Fourier变换是真正的有限维酉变换（差一个归一化常数）。连续频率、周期系数与离散变换之间的转换分别需要周期化、取样与截断，它们带来重叠、混叠与频率扩散。

Fourier表示用整体振荡的频率组织函数，每个坐标都依赖整个定义域上的值，局部的跳跃会影响所有高频坐标。第~\ref{chap:wavelet-transform}章改用尺度与位置同时组织局部细节；有限项截断与频带截断的逼近误差，以及系数衰减怎样转化为误差的速度，在第~\ref{chap:function-approximation}章统一研究。
